\HistoricalFoot{}{6*}
\begin{center}\bfseries 17.\end{center}
\noindent\hspace*{\parindent}Die Beantwortung der Frage nun nach den Teilern einer Funktion \((\mathrm{mod.}\ p)\), das Grundproblem der Lehre von den Kongruenzen, welche in Bezug auf einen Doppelmodul oder — in \textsc{Kronecker}scher Ausdrucksweise — in Bezug auf ein Modulsystem zweiter Stufe gedacht werden, gründet sich bei \textsc{Gauss} auf eine Reihe von Hilfssätzen. Einerseits lehrt er für eine Gleichung \(P(x)=0\) die andere Gleichung \(P_p(x)=0\) bilden, welche zu Wurzeln die \(p^{\mathrm{ten}}\) Potenzen von den Wurzeln der ersteren hat, und den (am 18. Aug. 1796, Nr. 26 T.) von ihm gefundenen Satz, dass \(P_p(x)\equiv P(x)\ (\mathrm{mod.}\ p)\) sei, wenn \(p\) eine Primzahl. Andererseits zeigt er, dass es für jede von \(x\) verschiedene Funktion \(P(x)\) einen kleinsten Exponenten \(\nu\) gibt, für welchen \(x^\nu-1\ (\mathrm{mod.}\ p)\) durch \(P(x)\) teilbar ist; und \(\nu\) ist ein Teiler von \(p^m-1\), falls \(P(x)\) Primfunktion vom Grade \(m\). Jede solche teilt also \((\mathrm{mod.}\ p)\) die Funktion \(x^{p^m-1}-1\); z. B. teilt jede von \(x\) verschiedene Primfunktion ersten Grades den Ausdruck \(x^{p-1}-1\), womit aufs Neue wieder der \textsc{Fermat}sche Satz bewiesen ist.

Aus diesen Ergebnissen folgt auf der einen Seite, dass \(x^{p^m-1}-1\ (\mathrm{mod.}\ p)\) dem Produkte aller inkongruenten von \(x\) verschiedenen Primfunktionen kon-
\newpage\SourceLeaf{51}\ArithmeticHead{44}{L}{BACHMANN, ÜBER GAUSS’ ZAHLENTHEORETISCHE ARBEITEN.}
\noindent gruent ist, deren Grade Teiler von \(m\) sind; auf der andern kann \(x^\nu-1\) nur solche primitive d. h. in keinem ähnlichen Ausdrucke geringeren Grades aufgehende Primteiler \((\mathrm{mod.}\ p)\) haben, deren Grad \(m\) den Exponenten bezeichnet, zu welchem \(p\ (\mathrm{mod.}\ \nu)\) gehört, derart, dass \(p^m\equiv1\ (\mathrm{mod.}\ \nu)\) (Nr. 30 T.). Auf solcher Grundlage kann nun für jede Funktion \(x^\nu-1\) nicht nur die Anzahl ihrer Primteiler eines bestimmten Grades ermittelt (art. 360, W. II, S. 230), sondern auch eine Methode gegeben werden, ihre (primitiven) Primteiler selbst zu finden. Dazu dient der am 30./31. August 1797 (Nr. 76, 77 T.) von \textsc{Gauss} erzielte Satz, dass jede ganze symmetrische Funktion der Grössen
\[
x,x^p,x^{p^2},\ldots,x^{p^{m-1}}
\]
in Bezug auf den Doppelmodul \(p,P(x)\), wenn \(P(x)\) eine Primfunktion \(m^{\mathrm{ten}}\) Grades bezeichnet, einer ganzen Zahl und die Koeffizienten von \(P(x)\) den elementaren symmetrischen Funktionen jener Potenzen kongruent sind; andererseits die Einteilung der Potenzen \(x^\alpha,x^\beta,\ldots\), wo \(\alpha,\beta,\ldots\) die zu \(\nu\) teilerfremden Zahlen \(<\nu\) bezeichnen, in Perioden, ähnlich den Perioden der Kreisteilungsgleichung, welche durch Kongruenzen bestimmt werden, die \((\mathrm{mod.}\ p)\) auflösbar sind, solange die Perioden noch aus solchen von der Gestalt
\[
x^k+x^kp+x^kp^2+\ldots+x^kp^{m-1}
\]
zusammengesetzt sind.

Schon viel früher hatte \textsc{Gauss} den Zusammenhang dieser Untersuchungen mit dem Fundamentaltheoreme erkannt (13. Aug. 1796, Nr. 23 T.); in der Tat führte ihn am 2. September 1796 (Nr. 30 T.) die Anwendung seiner Methoden auf den Fall, wo \(\nu\) eine ungerade Primzahl ist, zu zwei neuen Beweisen jenes Gesetzes, deren zweiter gewissermassen den umgekehrten Gang nimmt wie der erste. Zweifelsohne vor dem letztangegebenen Tag datiert eine Notiz in \textsc{Gauss}’ Nachlass (Ea 5, abgedruckt W. X 1, S. 114), in welcher er bezüglich der Gleichung, deren Wurzeln die \(e\) Perioden von \(f\) Gliedern \(p^{\mathrm{ter}}\) Einheitswurzeln sind, unter der Überschrift »der goldene Lehrsatz« den Satz unterstreicht: »diese Gleichung ist möglich für jeden Primmodulus \(=\sqrt[f]{1}\ (\mathrm{Mod.}\ p)\)« d. h. welcher zum Exponenten \(f\ (\mathrm{mod.}\ p)\) gehört, ein Satz, der als ein Theorema generale demonstrandum schon im \textsc{Leiste} bei Seite 108 (W. X 1, S. 115) aufgezeichnet ist. Die Quelle, aus der die genannten zwei Beweise des Rezi-
\newpage\SourceLeaf{52}\ArithmeticHead{45}{R}{DIE SPÄTEREN BEWEISE DES FUNDAMENTALTHEOREMS.}
\noindent prozitätsgesetzes fliessen, ist nur ein Spezialfall dieses allgemeinen Satzes, und so erklärt es sich, dass \textsc{Gauss} in seinem \textit{Tagebuche} am 27. Juni 1796 (Nr. 16 T.) den Ausdruck Theorema aureum auch für das Fundamentaltheorem verwendet.

Die A. R. schliesst mit einigen Sätzen ab, mit denen \textsc{Gauss} begonnen hat, die Theorie der Zerlegung der Funktionen auf den Fall auszudehnen, wo der Modul eine Primzahlpotenz oder allgemeiner eine beliebig zusammengesetzte Zahl ist (Nr. 77, 78 T. Ende August 1797), und diese Fälle auf den einfacheren eines Primzahlmoduls zurückzuführen (9. Sept. 1797, Nr. 79 T.). Hierin scheint er die Hilfsmittel gefunden zu haben, gewisse Schwierigkeiten, die jener einfachere Fall noch bot, zu beheben (s. art. 251, W. II, S. 209; art. 363 Ende, W. II, S. 232).
